\section{Tesina}
Investigación documental de carácter monográfico, con aportación al estado del arte. (Art. 73, REP).

\subsection{Requisitos de la opción}
\begin{itemize}
  \item Reporte de similitud iThenticate no mayor al 20 \%.
\end{itemize}
\subsection{Elegir la opción e iniciar}
En \emph{Mi trámite de licenciatura}, bajo \emph{¿Cómo quieres titularte?}, abre \emph{Con trabajo escrito} y elige \emph{Tesina}. La página de la opción explica qué necesitas, cómo será tu trámite y ofrece este manual para descargar, sin iniciar nada. Para empezar, acepta el aviso de privacidad y pulsa \boton{Iniciar mi trámite}.
\captura{tesina_pasante_01_entrada}{«¿Cómo quieres titularte?»: abres «Con trabajo escrito» y eliges la opción.}
\captura{tesina_pasante_02_informativa}{Página de la opción, antes de iniciar: qué necesitas, cómo será el trámite y el manual para descargar.}
\captura{tesina_pasante_03_iniciar}{Aceptas el aviso de privacidad y pulsas «Iniciar mi trámite».}
\subsection{Tu trámite, etapa por etapa}
Tu trámite tiene 8 etapas. Arriba siempre ves la tarjeta \emph{Tu siguiente paso} con lo que te toca ahora y una barra con las etapas: las terminadas, la que está en curso y las que vienen después.
\subsubsection*{Etapa 1. Tus datos}
Captura una sola vez tus datos; con ellos se llenan automáticamente todos los formatos del trámite.

Datos que capturas:
\begin{multicols}{2}\small
\begin{itemize}
  \item Nombre completo (con acentos, tal como aparecerá en el título)
  \item Número de cuenta
  \item Licenciatura
  \item Sistema
  \item Plan de estudios (clave y versión)
  \item Escuela incorporada a la UAEM donde cursaste la licenciatura
  \item Año de ingreso
  \item Año de egreso
  \item Teléfono
  \item Correo institucional
  \item Correo personal
  \item Ocupación actual
  \item Fecha de nacimiento
  \item Mi nombre o apellidos llevan acentos
\end{itemize}
\end{multicols}
\captura{tesina_pasante_04_siguiente}{Tu siguiente paso: Tus datos}
\captura{tesina_pasante_05_captura_datos}{Formulario de la etapa: capturas tus datos una sola vez.}
\subsubsection*{Etapa 2. Protocolo y aprobación del tema}
Captura la propuesta de tu trabajo. El sistema llena el formato oficial de solicitud y protocolo; escribe aquí todo el protocolo; la plataforma lo firma por ti, le pide la firma a tu asesor por correo y lo envía al responsable del Comité de Revisión de Protocolos de tu licenciatura. Si el comité te pide correcciones, las verás en tus Notificaciones: corrige aquí y vuelve a enviarlo.

Datos que capturas:
\begin{multicols}{2}\small
\begin{itemize}
  \item Título del trabajo
  \item Asesor (adscrito a la UAEM): nombre
  \item Asesor: correo electrónico
  \item Asesor: grado (Dr., Mtra., Ing.…)
  \item Asesor: género
  \item Coasesor (opcional): nombre
  \item Coasesor: correo electrónico
  \item Coasesor: grado
  \item Coasesor: género
  \item Área académica
  \item Planteamiento de la investigación documental o traducción crítica
  \item Justificación
  \item Objetivo general
  \item Metodología
  \item Referencias y/o fuentes de información
  \item Cronograma de actividades
\end{itemize}
\end{multicols}

La plataforma llena con tus datos los formatos de esta etapa; los firmas electrónicamente con un botón.

Documentos (PDF):
{\small\begin{longtable}{>{\raggedright\arraybackslash}p{0.47\linewidth}>{\raggedright\arraybackslash}p{0.45\linewidth}}
\toprule
\textbf{Documento} & \textbf{Observaciones}\\
\midrule
\endhead
Solicitud y protocolo firmados (pasante y asesor) & La plataforma lo genera. La plataforma lo genera con lo que escribiste; tu asesor lo firma por correo. Después lo revisa el comité de tu licenciatura.\\
Carta de presentación del pasante y asesor & Descárgala en blanco desde «Formatos en blanco». Solo para instituciones incorporadas.\\
\bottomrule
\end{longtable}}

\emph{Le toca al departamento:} Aprobación del protocolo y oficio de aprobación del tema. Cuando el comité aprueba el protocolo, la plataforma genera el oficio de aprobación del tema; lo firman la Jefatura de Evaluación Profesional (Vo. Bo.) y la Subdirección Académica.
\captura{tesina_pasante_06_siguiente}{Tu siguiente paso: Protocolo y aprobación del tema}
\captura{tesina_pasante_07_captura_protocolo}{Formulario de captura de la etapa.}
\captura{tesina_pasante_09_generar}{La plataforma llena el formato con tus datos y lo firma por ti con un botón.}
\captura{tesina_pasante_11_subir}{Subir un documento: eliges el PDF y lo envías; queda «En validación».}
\subsubsection*{Etapa 3. Desarrollo del trabajo escrito}
Tienes exclusividad sobre el tema durante dos años a partir de la aprobación; puedes renovarla con una prórroga hasta por un año, que incluye la sustentación (Art. 86, REP). Desarróllalo con tu asesor; cuando cuente con su visto bueno, confirma aquí para continuar.

Al terminar confirmas: «Concluí el trabajo escrito y mi asesor lo avala para revisión».
\captura{tesina_pasante_12_siguiente}{Tu siguiente paso: Desarrollo del trabajo escrito}
\captura{tesina_pasante_13_confirmar}{Confirmas que terminaste la etapa.}
\subsubsection*{Etapa 4. Solicitud de revisión del trabajo escrito (examen previo)}
Llena la solicitud de revisión del trabajo escrito y sube los documentos probatorios: el reporte de similitud de iThenticate (máximo 20 \%) y tu trabajo en PDF. Después elige la fecha y hora de tu examen previo y propone a tu sínodo; tu solicitud llega al Coordinador de Docencia de tu licenciatura.

Datos que capturas:
\begin{multicols}{2}\small
\begin{itemize}
  \item Título del trabajo escrito
\end{itemize}
\end{multicols}

La plataforma llena con tus datos los formatos de esta etapa; los firmas electrónicamente con un botón.

Documentos (PDF):
{\small\begin{longtable}{>{\raggedright\arraybackslash}p{0.47\linewidth}>{\raggedright\arraybackslash}p{0.45\linewidth}}
\toprule
\textbf{Documento} & \textbf{Observaciones}\\
\midrule
\endhead
Solicitud de revisión firmada (pasante y asesor) & La plataforma lo genera. La firmas tú y tu asesor da su Vo. Bo. desde su correo.\\
Reporte de porcentaje de similitud (iThenticate) & No mayor al 20 \%.\\
Trabajo escrito en PDF para revisión & Con copia del oficio de aprobación de tema al inicio.\\
Solicitud de prórroga del tema (pasante y asesor) & La plataforma lo genera. Tu protocolo se aprobó hace más del plazo de exclusividad (Art. 86, REP). Captura avance y motivos en «Solicitudes y formatos adicionales», firma aquí y tu asesor da su Vo. Bo.; al validarla, la plataforma genera el oficio de aprobación de prórroga que firma la Subdirección.\\
\bottomrule
\end{longtable}}

Eliges fecha y hora en los horarios disponibles y propones a tu sínodo o jurado; la solicitud llega al Coordinador de Docencia de tu licenciatura, que lo ratifica o modifica.

\emph{Le toca al departamento:} Fecha del examen previo y sínodo. Eliges fecha y hora y propones tu sínodo; el Coordinador de Docencia de tu licenciatura lo ratifica o modifica y firma desde el enlace de su correo. Después se generan los oficios individuales.
\captura{tesina_pasante_14_siguiente}{Tu siguiente paso: Solicitud de revisión del trabajo escrito (examen previo)}
\captura{tesina_pasante_15_captura_revision}{Formulario de captura de la etapa.}
\captura{tesina_pasante_17_generar}{La plataforma llena el formato con tus datos y lo firma por ti con un botón.}
\captura{tesina_pasante_21_agenda1}{Eliges fecha y hora del examen previo y propones a tu sínodo; la solicitud llega al Coordinador de Docencia.}
\subsubsection*{Etapa 5. Solicitud de sustentación y voto aprobatorio}
Cuando hayas atendido las correcciones del examen previo, genera el voto aprobatorio: la plataforma lo envía a tu asesor y a tus revisores para que lo firmen desde su correo. Con el voto firmado, firma tu solicitud de sustentación (tu asesor da su Vo. Bo.) para obtener la autorización de impresión y la asignación del jurado. Escribe aquí la estructura capitular: con ella la plataforma genera el oficio de estructura capitular que firma el Director.

Datos que capturas:
\begin{multicols}{2}\small
\begin{itemize}
  \item Estructura capitular (un capítulo por línea, sin numerar)
\end{itemize}
\end{multicols}

La plataforma llena con tus datos los formatos de esta etapa; los firmas electrónicamente con un botón.

Documentos (PDF):
{\small\begin{longtable}{>{\raggedright\arraybackslash}p{0.47\linewidth}>{\raggedright\arraybackslash}p{0.45\linewidth}}
\toprule
\textbf{Documento} & \textbf{Observaciones}\\
\midrule
\endhead
Voto aprobatorio (lo firman tu asesor y tus revisores) & La plataforma lo genera. La plataforma lo genera y pide las firmas por correo.\\
Solicitud de sustentación firmada (pasante y asesor) & La plataforma lo genera. Se habilita cuando el voto aprobatorio está firmado.\\
\bottomrule
\end{longtable}}

\emph{Le toca al departamento:} Autorización de impresión. 
\captura{tesina_pasante_22_siguiente}{Tu siguiente paso: Solicitud de sustentación y voto aprobatorio}
\captura{tesina_pasante_23_captura_sustentacion}{Formulario de captura de la etapa.}
\captura{tesina_pasante_25_generar}{La plataforma llena el formato con tus datos y lo firma por ti con un botón.}
\subsubsection*{Etapa 6. Documentación para la evaluación profesional}
Reúne y sube cada documento en PDF. Cuando todos estén validados, Apoyo a Titulación te cita para entregar originales, copias y ejemplares, y fijar el jurado.

La plataforma llena con tus datos los formatos de esta etapa; los firmas electrónicamente con un botón.

Documentos (PDF):
{\small\begin{longtable}{>{\raggedright\arraybackslash}p{0.47\linewidth}>{\raggedright\arraybackslash}p{0.45\linewidth}}
\toprule
\textbf{Documento} & \textbf{Observaciones}\\
\midrule
\endhead
Solicitud para presentar la evaluación profesional (trabajo escrito) & La plataforma lo genera. La plataforma lo genera con tus datos; lo firmas electrónicamente y la Subdirección lo autoriza.\\
Trabajo escrito final en PDF & Con oficio de autorización de impresión, resumen y carta de autorización para el Repositorio Institucional.\\
Reporte de porcentaje de similitud (iThenticate) & No mayor al 20 \%.\\
Certificado de estudios de licenciatura & Descargado del Portal de alumnos.\\
Certificado de servicio social & \\
Certificado de estudios de bachillerato & \\
Constancia de autenticación de bachillerato & Se tramita en la escuela donde cursaste el bachillerato. Solo si tus estudios de bachillerato no fueron en instituciones dependientes o incorporadas a la UAEM.\\
Acta de nacimiento & Copia certificada vigente (del año en curso).\\
CURP & Debe mostrar la leyenda «Validada por el Registro Civil».\\
Constancia o certificado de aprobación del idioma inglés & Solo planes 98 en adelante, si en tu certificado no aparecen Inglés C1 e Inglés C2.\\
Constancia de no adeudo a la Biblioteca de la Facultad & El recibo de pago se tramita en Control Escolar de la Facultad; la constancia, en la Biblioteca.\\
Constancia de no adeudo a la Biblioteca Central & \\
Contra-recibo de la donación de libros a la Biblioteca de la Facultad & \\
Constancia de no adeudo a la Dirección de Control Escolar de la UAEM & Se solicita en el Portal de alumnos: controlescolar.uaemex.mx.\\
Recibo de no adeudo a la UAEM & Se solicita en el Portal de alumnos de Control Escolar.\\
Formato de datos personales (para el título) & La plataforma lo genera. La plataforma lo llena con tus datos; lo firmas electrónicamente.\\
Comprobante del sistema de Seguimiento de Egresados & Se obtiene en alumniversitario.uaemex.mx/1dae.php.\\
Dictamen de autorización de los HH. Consejos & La plataforma lo genera. Se solicita en la Subdirección Académica con el formato del Art. 5. Solo si excediste el plazo de 10 años para presentar la evaluación profesional (Art. 5, REP).\\
Carta de autorización o de excepción para publicar en el Repositorio Institucional & Entrega la que corresponda; descárgala en «Formatos en blanco».\\
Recibo de pago: derechos de evaluación profesional (\$3,024) & Se solicita en Control Escolar de la Facultad.\\
Recibo de pago: expedición de título (\$1,187) & \\
\bottomrule
\end{longtable}}
\captura{tesina_pasante_28_siguiente}{Tu siguiente paso: Solicitud para presentar la evaluación profesional (trabajo escrito)}
\captura{tesina_pasante_29_generar}{La plataforma llena el formato con tus datos y lo firma por ti con un botón.}
\subsubsection*{Etapa 7. Solicitud del examen definitivo}
Elige la fecha y hora de tu examen definitivo y propone a tu jurado; tu solicitud llega al Coordinador de Docencia de tu licenciatura, que lo ratifica o modifica.

Eliges fecha y hora en los horarios disponibles y propones a tu sínodo o jurado; la solicitud llega al Coordinador de Docencia de tu licenciatura, que lo ratifica o modifica.

\emph{Le toca al departamento:} Fecha del examen definitivo y jurado. Eliges fecha y hora y propones tu jurado; el Coordinador de Docencia lo ratifica o modifica y firma desde el enlace de su correo. Después se generan los oficios individuales.
\captura{tesina_pasante_33_siguiente}{Tu siguiente paso: Solicitud del examen definitivo}
\captura{tesina_pasante_34_agenda2}{Eliges fecha y hora del examen definitivo y propones a tu sínodo; la solicitud llega al Coordinador de Docencia.}
\subsubsection*{Etapa 8. Resultado de la evaluación}
Al concluir, el departamento registra el resultado. Después podrás consultar cómo descargar tu título y cédula profesional.

\emph{Le toca al departamento:} Registrar resultado. 
\captura{tesina_pasante_35_siguiente}{Tu siguiente paso: Resultado de la evaluación}
\subsection{Si te devuelven un documento}
Cuando Titulación devuelve un documento, arriba de tu trámite aparece \emph{Falta corregir un documento} con la observación, y tu siguiente paso te lleva a corregirlo. Sube el PDF corregido desde el mismo renglón.
\captura{tesina_pasante_19_devuelto}{Un documento devuelto: arriba se muestra qué corregir y la observación de Titulación.}
\captura{tesina_pasante_20_resubir}{Corriges el documento devuelto: lees la observación, eliges el PDF corregido y lo vuelves a enviar.}
\subsection{Al terminar}
Cuando se registra el resultado, tu trámite queda concluido. Titulación verifica tu expediente y lo envía a Servicios Escolares para la expedición del título. En \emph{Estado de mi trámite} ves en qué va.
\captura{tesina_pasante_37_concluido}{Trámite concluido: el resultado queda registrado.}
\captura{tesina_pasante_38_estado}{«Estado de mi trámite»: en qué va y qué sigue, por etapas.}
